
\documentclass[11pt,a4paper]{ctexart}
\usepackage{amsmath,amssymb,bm}
\usepackage{graphicx}
\usepackage{booktabs}
\usepackage{tabularx}
\usepackage[table]{xcolor}
\usepackage{listings}
\usepackage{caption}
\usepackage{geometry}
\usepackage{enumitem}
\usepackage[hidelinks]{hyperref}

\geometry{left=2.5cm,right=2.5cm,top=2.4cm,bottom=2.6cm}
\linespread{1.42}
\setlength{\parindent}{2em}
\setlength{\parskip}{0.25em}

\captionsetup{font=small,labelfont=bf,skip=5pt}
\captionsetup[table]{position=bottom,skip=4pt}

\definecolor{hdark}{HTML}{12304F}
\definecolor{hmid}{HTML}{1F4E79}
\definecolor{hlight}{HTML}{2F5D80}
\definecolor{thd}{HTML}{2F5D80}
\definecolor{tal}{HTML}{F2F5F8}

\newcommand{\hOne}[1]{\par\vspace{0.5em}%
  {\color{hdark}\Large\bfseries #1}\par\vspace{0.1em}%
  \noindent{\color{hdark}\rule{\linewidth}{0.9pt}}\par\vspace{0.35em}}
\newcommand{\hTwo}[1]{\par\vspace{0.8em}{\color{hmid}\large\bfseries #1}\par\vspace{0.25em}}
\newcommand{\hThree}[1]{\par\vspace{0.6em}{\color{hlight}\normalsize\bfseries #1}\par\vspace{0.15em}}

\lstset{
  basicstyle=\ttfamily\footnotesize,
  breaklines=true,
  breakatwhitespace=false,
  frame=single,
  framesep=5pt,
  rulecolor=\color{gray!45},
  backgroundcolor=\color{gray!6},
  columns=fullflexible,
  keepspaces=true,
  showstringspaces=false,
  xleftmargin=8pt,
  xrightmargin=8pt,
  aboveskip=0.7em,
  belowskip=0.7em,
}

\pagestyle{empty}
\setlength{\parindent}{2em}
\setlength{\parskip}{0.25em}
\title{复现代码与结果说明}
\author{MotrixLab 线上实习}
\begin{document}
\begin{center}{\color{hdark}\Huge\bfseries 复现代码与结果说明}\end{center}

\begin{center}{\color{hmid}\large 《第三周周报》配套材料 $\cdot$ 实现要点与实测数据}\end{center}

\vspace{0.35cm}

\noindent \textbf{实习项目：}MotrixLab 线上实习（第一月）　　\textbf{本周课次：}第 3 周　　\textbf{代码位置：}\textbf{code/}（8 个脚本 + 4 个共享模块）\par

\vspace{0.30cm}

\hOne{一、概述}

本文是《第三周周报》的配套材料。周报只讲理论（概念、推导与判据），本文负责说明三件事：\textbf{写了哪些代码、每个文件的实现要点、以及跑出来的实际结果}。\par

所有脚本都遵循同一套约定，以便“结论可核对、过程可复现”：\par

\begin{itemize}[leftmargin=1.6em,itemsep=0.15em,topsep=0.25em]

\item 每个脚本都可以\textbf{独立运行}，内置断言全部通过后才会打印 \textbf{all assertions passed}；断言失败即退出码非零，可直接用于 CI。

\item 关键数学都配了\textbf{独立的交叉验证源}，避免“用同一个错误实现验证自己”的假阴性：正运动学对解析闭式解、解析雅可比对有限差分、SLERP 对解析转角、GAE 对 $\lambda$ = 0 与 $\lambda$ = 1 两个退化端点。

\item 图落盘到 \textbf{figures/}，数值结果落盘到 \textbf{results/}；强化学习部分的脚本还会把指标写成 JSON，供汇总脚本读取。

\item 随机性全部显式播种（torch、numpy 与环境重置），因此同一脚本重复运行结果完全一致。

\end{itemize}

\textbf{依赖：}numpy、scipy、matplotlib、gymnasium、torch（CPU 即可，无需 GPU）；报告生成部分额外需要 tectonic（LaTeX 引擎）。\par

\clearpage

\hOne{二、代码模块说明}

\hTwo{2.1 共享模块}

\begin{table}[htbp]
\centering
\begin{tabularx}{\textwidth}{>{\hsize=0.5571\hsize\raggedright\arraybackslash}X>{\hsize=1.8000\hsize\raggedright\arraybackslash}X>{\hsize=0.6429\hsize\raggedright\arraybackslash}X}
\toprule
\rowcolor{thd}\textcolor{white}{\textbf{模块}} & \textcolor{white}{\textbf{职责}} & \textcolor{white}{\textbf{被谁使用}} \\
\midrule
\textbf{se3.py} & SE(3)/SO(3) 基础工具：旋转矩阵、欧拉角、四元数、旋转向量互转，万向锁处理，SLERP，齐次矩阵构造/求逆/作用 & 01、02、03 \\
\rowcolor{tal}
\textbf{traj.py} & 梯形速度剖面（LSPB）：profile 生成、逐时刻求值、固定时长下的最优加速时间搜索 & 03、04 \\
\textbf{rl\_common.py} & CartPole 环境封装、状态离散化、统一评估协议、随机基线、结果 JSON 落盘 & 05、06、07、08 \\
\rowcolor{tal}
\textbf{mpl\_style.py} & matplotlib 中文字体探测与设置（含无中文字体时的降级提示） & 全部绘图脚本 \\
\bottomrule
\end{tabularx}
\caption{四个共享模块的职责划分。}
\end{table}

之所以把这些抽出来，是因为它们被多处复用且容易写错：\textbf{se3.py} 里所有函数都在 docstring 里写明了约定（四元数一律 xyzw、欧拉角一律 ZYX、$T_{ij}$ 表示 j 系在 i 系中的位姿），目标是“约定只在一处定义，中间层永不混用”；\textbf{traj.py} 被 03 和 04 共用，避免同一套速度剖面逻辑写两遍而产生不一致。\par

\hTwo{2.2 位姿变换（01\_\allowbreak pose\_\allowbreak transform.\allowbreak py）}

覆盖六个方面：旋转矩阵基本性质、欧拉角与旋转矩阵往返、旋转矩阵与四元数往返、万向锁、齐次矩阵链式法则、SLERP 匀速性。核心设计是把“底层数学”全部放在 se3.py，本脚本只负责\textbf{验证约定 + 出图}，因此断言写得很直接。\par

两处值得说明的实现选择：\par

\begin{itemize}[leftmargin=1.6em,itemsep=0.15em,topsep=0.25em]

\item \textbf{万向锁用显式分支处理。}当 |sin(pitch)| 接近 1 时，把 roll 置零、把等效转角全部并入 yaw，并断言“重构出的姿态与原姿态严格等价”。这样即使用户传入奇异附近的欧拉角，也不会产生 NaN 或跳变。

\item \textbf{解析求逆与数值求逆对照。}脚本同时用闭式解 T$^{-}$$^{1}$ = [[R$^{\top}$, $-$R$^{\top}$p], [0, 1]] 和 np.linalg.inv 求逆并断言两者一致，既验证了公式，也说明生产代码里应当用前者。

\end{itemize}

\hTwo{2.3 正运动学与逆解（02\_\allowbreak forward\_\allowbreak kinematics.\allowbreak py）}

用 UR5 的标准 DH 参数搭 6 轴链，验证四件事：DH 变换本身、链式累乘、几何雅可比、数值逆运动学。\par

\textbf{雅可比用几何法构造}（线速度列 $J_{v}$ = $z_{i}$ $\times$ ($p_{\mathrm{end}}$ $-$ $p_{i}$)，角速度列 $J_{w}$ = $z_{i}$），再用有限差分法独立算一遍作对照。姿态部分的差分用旋转矩阵的反对称部分提取旋转向量，容差放宽到 1e-5（因为依赖差分步长）。\par

\textbf{逆解用阻尼最小二乘}（Levenberg–Marquardt 风格）迭代：\par

\begin{equation}\Delta q=J^{\top}\left(JJ^{\top}+\lambda^{2}I\right)^{-1}e\end{equation}

阻尼项 $\lambda$ 的作用是在奇异位形附近保持解有界；脚本另外做了单步限幅（避免奇异附近跳出一大步）和关节限位裁剪。为了说明阻尼的必要性，脚本专门构造了一个腕部奇异位形（q$_{5}$ = 0，关节 4 与关节 6 轴线共线），断言此时最小奇异值趋于零，并对比裸伪逆与阻尼解的关节速度范数。\par

\hTwo{2.4 轨迹规划（03\_\allowbreak trajectory\_\allowbreak planning.\allowbreak py）}

实现了三类剖面与一个判据函数：\par

\begin{itemize}[leftmargin=1.6em,itemsep=0.15em,topsep=0.25em]

\item \textbf{三次多项式：}闭式解系数（a$_{0}$ = q$_{0}$、a$_{1}$ = 0、a$_{2}$ = 3$\Delta$q/T$^{2}$、a$_{3}$ = $-$2$\Delta$q/T$^{3}$），两端速度为零。

\item \textbf{五次多项式：}六个边界条件（两端的位置、速度、加速度）解 6$\times$6 线性方程组。

\item \textbf{LSPB 梯形剖面：}由 traj.py 提供，自动判断是梯形还是三角形。

\item \textbf{plan\_\allowbreak with\_\allowbreak fixed\_\allowbreak duration：}给定总时长，解析地求最优加速时间（$t_{a}$* = min($v_{\mathrm{max}}$/$a_{\mathrm{max}}$, T/2)），并给出可行性结论。

\end{itemize}

\textbf{多项式求值不使用数值差分}，而是对系数逐阶求导得到速度与加速度，这样峰值是精确的、不受差分步长影响。这一细节很重要：如果用差分去检查峰值是否超限，差分误差会让边界情况误判。\par

脚本还实现了笛卡尔直线 + 姿态 SLERP，并用“关节空间线性插值 + 正运动学”造出对比弧线，用来量化“关节空间插值末端不走直线”这个结论。\par

\hTwo{2.5 采样法路径规划（04\_\allowbreak rrt\_\allowbreak planner.\allowbreak py）}

在 1 m $\times$ 1 m 的平面里放了三个圆形障碍和一个矩形障碍，从 (0.08, 0.08) 走到 (0.92, 0.92)。完整走一遍规划流水线：\par

\begin{lstlisting}
碰撞检测 -> RRT / RRT* 采样 -> 捷径平滑 -> 按 (vmax, amax) 时间参数化
\end{lstlisting}

碰撞检测用\textbf{线段密集采样}（步长 5 mm）而不是解析求交。这是一个有意的取舍：解析求交要多写几十行且各有边界情况，而采样步长比最小障碍尺寸（90 mm）小一个数量级，不会穿透。代码里用注释标明这个上限以及升级路径。\par

RRT* 相对 RRT 只多了两处：\textbf{choose parent}（邻域内选累计代价最小的父节点）与 \textbf{rewire}（邻域内其他节点改认新节点，若这样更省）。时间参数化沿路径\textbf{弧长}套用 traj.py 的梯形剖面，先算累计弧长、再按弧长插值出位置与切向速度。\par

\hTwo{2.6 表格 Q-learning 与 DQN（05\_\allowbreak qlearning\_\allowbreak dqn.\allowbreak py）}

\textbf{表格 Q-learning：}CartPole 的观测是 4 维连续量且部分无界，所以先裁剪到有限区间再分箱。Q 表用 4 维数组 Q[x, $\dot{x}$, $\theta$, $\dot{\theta}$, a] 存储，正好可以用离散化返回的下标直接索引，避免手工做扁平化换算。\par

\textbf{DQN：}实现了四个关键机制，每一个都对应 deadly triad 的一个病：\par

\begin{itemize}[leftmargin=1.6em,itemsep=0.15em,topsep=0.25em]

\item \textbf{经验回放}——用预分配的 numpy 数组做环形缓冲，采样时随机取小批量，打散时间相关性。

\item \textbf{目标网络}——每 200 步同步一次参数，稳住自举目标。

\item \textbf{Double DQN}——在线网络选动作、目标网络估值，抑制 max 的高估。

\item \textbf{Huber 损失}——防止时序差分误差的尖峰把梯度带飞。

\end{itemize}

另外脚本按固定间隔做贪心评估并\textbf{保留最优快照}，理由是 DQN 训练中途会崩（见 4.5 的实测曲线），只交最后一个网络是不可靠的。这也是业界默认的模型选择做法。\par

\hTwo{2.7 策略梯度（06\_\allowbreak reinforce\_\allowbreak actor\_\allowbreak critic.\allowbreak py）}

三种变体共用一个训练循环，\textbf{只换优势函数}，以保证对比公平（网络结构、随机种子、episode 预算完全相同）：\par

\begin{table}[htbp]
\centering
\begin{tabularx}{\textwidth}{>{\hsize=0.7286\hsize\raggedright\arraybackslash}X>{\hsize=1.3714\hsize\raggedright\arraybackslash}X>{\hsize=0.9000\hsize\raggedright\arraybackslash}X}
\toprule
\rowcolor{thd}\textcolor{white}{\textbf{方法}} & \textcolor{white}{\textbf{优势函数}} & \textcolor{white}{\textbf{评论家是否训练}} \\
\midrule
REINFORCE & $A_{t}$ = $G_{t}$（整条轨迹的折扣回报） & 不训练（结构上保留但不用） \\
\rowcolor{tal}
REINFORCE + 基线 & $A_{t}$ = $G_{t}$ $-$ V($s_{t}$) & 用蒙特卡洛回报拟合 \\
A2C & $A_{t}$ = $r_{t}$ + $\gamma$V(s\_{t+1}) $-$ V($s_{t}$) & 用自举目标训练 \\
\bottomrule
\end{tabularx}
\caption{三种策略梯度变体的唯一差别。}
\end{table}

两处工程要点：一是采样期间\textbf{不开梯度、不建计算图}，只记录 states/actions/rewards，episode 结束后对整条轨迹做一次批量前向算 log 概率与熵 —— 逐帧建图会把 500 步的计算图全留着，速度差好几倍；二是 A2C 的自举\textbf{必须在真正终止处切断}（截断是时间上限、不是任务结束，仍要自举），脚本为此单独写了断言。\par

为了量化“基线降方差”这件事，脚本做了一个受控实验：\textbf{固定一个策略、采同一批 60 条轨迹，只换优势函数的形式}，比较信用分配信号的标准差。之所以要固定策略，是因为直接拿训练过程中的回报离散度去比会混进混杂因素（某个方法收敛差、episode 短，回报本来就集中），结论会被误导。\par

\hTwo{2.8 PPO 与裁剪消融（07\_ppo.py）}

PPO 部分实现了：Actor-Critic 共享躯干网络、GAE($\lambda$) 优势估计、重要性采样 + 裁剪目标、多 epoch 小批量更新、优势归一化、梯度裁剪、学习率线性退火。\par

脚本的核心产出是一次\textbf{严格消融}：其他所有超参完全相同，只把 $\varepsilon$ 设为无穷大（等价于不裁剪），对照两者的近似 KL、裁剪生效比例与学习曲线。为了让消融能说明问题，这里\textbf{故意没有开启}“KL 超限提前结束 epoch”这道保险 ——否则它会先于裁剪生效，消融就看不出裁剪本身的作用了。\par

脚本还用四条断言把裁剪目标的性质逐项钉死：\par

\begin{itemize}[leftmargin=1.6em,itemsep=0.15em,topsep=0.25em]

\item ratio 在 [1$-$$\varepsilon$, 1+$\varepsilon$] 内时，梯度等于优势（裁剪不生效）；

\item 正优势且 ratio 超过上界时，梯度为 0（策略不能再往上推）；

\item 正优势且 ratio 低于下界时，梯度仍为优势（鼓励往回走）；

\item 负优势情形对称：ratio 过小时梯度为 0；ratio $\equiv$ 1 时目标退化为普通策略梯度。

\end{itemize}

\hTwo{2.9 汇总与 A3C 实验（08\_\allowbreak algorithm\_\allowbreak summary.\allowbreak py）}

读取 05/06/07 的结果文件，汇总成横向对比表；另外实现了一个 A3C 并行采样的方差实验：\textbf{固定一个策略}，让每个“worker”跑一条轨迹，统计“平均 B 条轨迹的梯度估计”的标准差，再对数据做 log-log 幂律拟合，看指数是否落在理论值 0.5 附近。\par

这个实验对\textbf{可复现性特别敏感}：每条轨迹的初始状态都必须显式播种，否则每次跑出来的数字都不一样。这一条在 5.2 里作为踩坑记录单独列出。\par

\clearpage

\hOne{三、运行方式与耗时}

\begin{lstlisting}
cd code
python 01_pose_transform.py          # 第 1 周：坐标系与位姿变换
python 02_forward_kinematics.py
python 03_trajectory_planning.py     # 第 2 周：轨迹规划
python 04_rrt_planner.py
python 05_qlearning_dqn.py           # 第 3 周：深度强化学习
python 06_reinforce_actor_critic.py
python 07_ppo.py                     # 依赖 06 的结果文件
python 08_algorithm_summary.py       # 依赖 05/06/07 的结果文件
\end{lstlisting}

\textbf{注意：05–08 需要按顺序运行}，因为 07 会读取 06 写下的 REINFORCE 曲线，08 会读取 05/06/07 三份结果文件。01–04 之间没有依赖，可以任意顺序运行。\par

\begin{table}[htbp]
\centering
\begin{tabularx}{\textwidth}{>{\hsize=1.0500\hsize\raggedright\arraybackslash}X>{\hsize=1.3500\hsize\raggedright\arraybackslash}X>{\hsize=1.0500\hsize\raggedright\arraybackslash}X>{\hsize=0.5500\hsize\raggedright\arraybackslash}X}
\toprule
\rowcolor{thd}\textcolor{white}{\textbf{脚本}} & \textcolor{white}{\textbf{对应内容}} & \textcolor{white}{\textbf{产出}} & \textcolor{white}{\textbf{实测耗时}} \\
\midrule
01\_\allowbreak pose\_\allowbreak transform.\allowbreak py & 齐次矩阵、欧拉角、四元数、万向锁、SLERP & 图 1 & 1.4 s \\
\rowcolor{tal}
02\_\allowbreak forward\_\allowbreak kinematics.\allowbreak py & 6 轴 DH 正运动学、雅可比、阻尼最小二乘 IK & 图 2 & 1.2 s \\
03\_\allowbreak trajectory\_\allowbreak planning.\allowbreak py & 三次 / 五次多项式、LSPB、笛卡尔直线 + SLERP & 图 3、trajectory.csv & 1.2 s \\
\rowcolor{tal}
04\_\allowbreak rrt\_\allowbreak planner.\allowbreak py & RRT / RRT*、捷径平滑、时间参数化 & 图 4、rrt\_path.csv & 2.5 s \\
05\_\allowbreak qlearning\_\allowbreak dqn.\allowbreak py & 表格 Q-learning vs DQN & 图 5、rl\_\allowbreak qlearning\_\allowbreak dqn.\allowbreak json & 约 2.5 min \\
\rowcolor{tal}
06\_\allowbreak reinforce\_\allowbreak actor\_\allowbreak critic.\allowbreak py & REINFORCE / +基线 / A2C & 图 6、rl\_\allowbreak policy\_\allowbreak gradient.\allowbreak json & 约 4 min \\
07\_ppo.py & PPO + GAE + 裁剪消融 & 图 7、rl\_ppo.json & 约 4.8 min \\
\rowcolor{tal}
08\_\allowbreak algorithm\_\allowbreak summary.\allowbreak py & 横向汇总 + A3C 并行方差实验 & 图 8、rl\_summary.json & 20 s \\
\bottomrule
\end{tabularx}
\caption{脚本清单与实测耗时（Windows / CPU 单机，总耗时约 12 分钟）。01–04 为确定性数值验证，秒级完成；05–08 是训练类任务，耗时集中在环境交互。}
\end{table}

首次运行前若缺少依赖，可用：\par

\begin{lstlisting}
pip install numpy scipy matplotlib gymnasium torch
\end{lstlisting}

报告与文档本身由 \textbf{\_tools/build\_pdf.py} 用 LaTeX（tectonic）编译生成。它会把 content\_\allowbreak theory.\allowbreak py 与 content\_code.py 两份内容文件转成 .tex，再分别产出《第三周周报》与《复现代码与结果说明》；图中的实验数据同样是构建时从 \textbf{results/*.json} 读取的，所以重跑实验后再生成，数据会自动同步，不会出现“文档写的和代码跑的不一致”。\par

\clearpage

\hOne{四、实验结果}

\hTwo{4.1 位姿变换}

脚本 01 的 6 组断言全部通过。图 1 的三个子图分别对应：两连杆臂的右乘链式（(a)）、万向锁在 pitch $\rightarrow$ 90° 时的数值崩溃（(b)）、以及 SLERP 相对欧拉角线性插值的匀速性差异（(c)）。\par

\begin{figure}[htbp]
\centering
\includegraphics[width=\linewidth]{D:/github/Xbotics-Learning-Recording/Xbotics-Learning-Recording/Week3/figures/fig_pose_transform.png}
\caption{位姿变换的三个关键现象。(b) 中 pitch 逼近 90° 时欧拉角逆解误差从 1e-16 量级跳到 1e-1 量级，这就是万向锁；(c) 中 SLERP 的转角随时间是严格直线，而欧拉角分量线性插值不是匀速转动。}
\end{figure}

\begin{table}[htbp]
\centering
\begin{tabularx}{\textwidth}{>{\hsize=0.6375\hsize\raggedright\arraybackslash}X>{\hsize=1.1250\hsize\raggedright\arraybackslash}X>{\hsize=1.2375\hsize\raggedright\arraybackslash}X}
\toprule
\rowcolor{thd}\textcolor{white}{\textbf{验证项}} & \textcolor{white}{\textbf{方法}} & \textcolor{white}{\textbf{实测结果}} \\
\midrule
旋转矩阵性质 & 正交性 + 行列式 = +1 & 误差 $<$ 1e-12，通过 \\
\rowcolor{tal}
欧拉角 $\leftrightarrow$ 旋转矩阵 & ZYX 约定往返 & 误差 $<$ 1e-10，通过 \\
旋转矩阵 $\leftrightarrow$ 四元数 & xyzw 约定往返 & 误差 $<$ 1e-10，通过 \\
\rowcolor{tal}
万向锁处理 & pitch = ±90° 时置 roll = 0 & 重构姿态严格等价，通过 \\
解析求逆 vs 数值求逆 & 闭式解对照 np.linalg.inv & 误差 $<$ 1e-12，通过 \\
\rowcolor{tal}
链式右乘 & T$_{0}$$_{3}$ = T$_{0}$$_{1}$T$_{1}$$_{2}$T$_{2}$$_{3}$ 与逐步变换对照 & 误差 $<$ 1e-12，通过 \\
两连杆正运动学 & 齐次矩阵链 vs 解析闭式解 & 误差 $<$ 1e-12，通过 \\
\rowcolor{tal}
SLERP 匀速性 & 等间隔采样的相邻转角 & 相等，误差 $<$ 1e-9，通过 \\
\bottomrule
\end{tabularx}
\caption{位姿变换模块的验证项与实测误差（脚本 01）。}
\end{table}

一个具体算例：B 系相对 A 系绕 z 轴转 45°、原点位于 (0.5, 0, 0.2)。把 B 系下的点 (1, 0, 0) 换算到 A 系，程序输出 (1.2071, 0.7071, 0.2)；手算 (cos45°, sin45°, 0) + (0.5, 0, 0.2) 得到同样结果，数值与解析两侧互相印证。\par

\hTwo{4.2 正运动学、雅可比与逆解}

取关节角 q = (0.3, $-$0.8, 1.1, $-$0.5, 1.4, 0.2) rad，程序算出的末端位姿为 [x, y, z, roll, pitch, yaw] = ($-$0.6984, $-$0.3449, 0.2015, 1.3711, $-$0.1623, $-$1.0639)，四元数 (0.512, $-$0.374, $-$0.347, 0.691)；该位形下雅可比奇异值为 (1.9176, 1.5522, 0.9852, 0.4585, 0.4450, 0.1688)，可操作度 0.1010。\par

\begin{figure}[htbp]
\centering
\includegraphics[width=\linewidth]{D:/github/Xbotics-Learning-Recording/Xbotics-Learning-Recording/Week3/figures/fig_forward_kinematics.png}
\caption{6 轴臂的正运动学与逆解。(a) 基座系 xz 投影下初始位形与 IK 解的对比，红星为目标点；(b) 解析几何雅可比与有限差分雅可比逐项吻合；(c) 阻尼最小二乘 IK 的误差单调下降，约 12 次迭代收敛到 1e-12 以下。}
\end{figure}

\begin{table}[htbp]
\centering
\begin{tabularx}{\textwidth}{>{\hsize=0.6000\hsize\raggedright\arraybackslash}X>{\hsize=1.3125\hsize\raggedright\arraybackslash}X>{\hsize=1.0875\hsize\raggedright\arraybackslash}X}
\toprule
\rowcolor{thd}\textcolor{white}{\textbf{检查项}} & \textcolor{white}{\textbf{做法}} & \textcolor{white}{\textbf{实测结果}} \\
\midrule
DH 变换正确性 & 与教科书展开式逐项对照 & 误差 $<$ 1e-12 \\
\rowcolor{tal}
链式累乘 & 不借助 dh\_transform，按 RzTzTxRx 顺序手写相乘 & 逐级一致，误差 $<$ 1e-12 \\
解析雅可比 & 几何法 vs 有限差分 & 位置误差 $<$ 1e-6，姿态 $<$ 1e-5 \\
\rowcolor{tal}
雅可比速度映射 & J$\cdot$dq 与末端位姿数值微分对照 & 误差 $<$ 1e-5 \\
数值逆运动学 & 阻尼最小二乘，容差 1e-10 & 约 12 次迭代收敛，位姿误差 $<$ 1e-7 \\
\rowcolor{tal}
奇异位形 & 腕部对齐（q$_{5}$ = 0）时最小奇异值 & $<$ 1e-6，并按预期秩亏 \\
奇异处阻尼必要性 & 对比裸伪逆与阻尼解 & 阻尼解的关节速度范数显著更小 \\
\bottomrule
\end{tabularx}
\caption{正运动学 / 雅可比 / 逆解的验证结果（脚本 02）。}
\end{table}

奇异位形这一项值得展开说：在 q$_{5}$ = 0 时关节 4 与关节 6 的轴线共线，雅可比秩亏，裸伪逆会解出范数极大的关节速度，真机上会直接触发关节限速乃至撞机。脚本用阻尼形式 dq = J$^{\top}$(JJ$^{\top}$ + $\lambda$$^{2}$I)$^{-}$$^{1}$e，并断言阻尼解的速度范数严格小于裸伪逆 ——这就是工业控制器里必开阻尼的原因。\par

\hTwo{4.3 轨迹规划}

脚本 03 的实测结果验证了周报 3.1 与 3.3 的推导：对 $\Delta$q = 1 rad、T = 1 s 的三次多项式，程序算出的峰值速度与峰值加速度为 1.500 rad/s 与 6.000 rad/s$^{2}$，与闭式解 1.5$\Delta$q/T、6$\Delta$q/T$^{2}$ 完全吻合。而设定的上限是 $v_{\mathrm{max}}$ = 0.6 rad/s、$a_{\mathrm{max}}$ = 2.0 rad/s$^{2}$，也就是说这条曲线\textbf{两项都超限}（分别超 2.5 倍与 3 倍）——这正是“必须显式限制速度加速度”的算术依据。\par

\begin{figure}[htbp]
\centering
\includegraphics[width=\linewidth]{D:/github/Xbotics-Learning-Recording/Xbotics-Learning-Recording/Week3/figures/fig_trajectory.png}
\caption{轨迹规划。(a) 三次 / 五次多项式与 LSPB 的位置曲线；(b) 速度、加速度与限制线 —— 三次多项式的两条曲线都冲出了虚线限制，LSPB 恰好压在线上；(c) 关节空间插值走弧线，笛卡尔直线插值走直线。}
\end{figure}

固定时长规划器的实测输出与表 1（周报）的算例一致：行程 0.8 rad 的自然时长为 1.633 s，最优峰值速度为 0.471 rad/s、峰值加速度 1.569 rad/s$^{2}$，2 s 内可行；行程 1.2 rad 的自然时长 2.300 s，最优峰值 0.706 rad/s 与 2.353 rad/s$^{2}$，两个限制同时超 18\%，2 s 内不可行。脚本还断言了解析最优解不劣于 2001 点网格搜索，以及“可行 $\Leftrightarrow$ 自然时长 $\leq$ T”这条等价判据。\par

笛卡尔与关节空间的对比也用数值确认了：笛卡尔直线插值的路径长度与两点直线距离完全相等（0.707107 m，误差 $<$ 1e-9），而同样的起终点在关节空间线性插值时，末端走出的弧线比弦长多 1\% 以上。\par

\hTwo{4.4 采样法路径规划}

同一随机种子下：RRT 得到 90 个节点、路径长 1.4572 m；RRT* 同样 90 个节点、路径长 1.3191 m，缩短约 9.5\%；再用随机捷径平滑压到 4 个途经点、1.2313 m，比 RRT 原始解短 15.5\%。\par

\begin{figure}[htbp]
\centering
\includegraphics[width=\linewidth]{D:/github/Xbotics-Learning-Recording/Xbotics-Learning-Recording/Week3/figures/fig_rrt.png}
\caption{RRT 与 RRT* 规划结果。(a)(b) 采样树与最终路径，蓝色细线是搜索树，红线是路径；(c) 捷径平滑后再按弧长做时间参数化得到的末端轨迹，颜色表示瞬时速度，两端为加减速段，中段达到 $v_{\mathrm{max}}$。}
\end{figure}

时间参数化取 $v_{\mathrm{max}}$ = 0.5 m/s、$a_{\mathrm{max}}$ = 1.5 m/s$^{2}$，得到总时长 2.796 s，峰值速度 0.500 m/s、峰值加速度 1.500 m/s$^{2}$，两项都恰好压在限制上且不越界；脚本同时断言了整条时间轨迹上的采样点都不在障碍物里。这就是“路径–速度解耦”的完整落地。\par

\hTwo{4.5 Q-learning 与 DQN}

评价协议统一为：贪心策略跑 50 个 episode 取平均回报。同时给出随机策略与恒定动作作为参照。\par

\begin{figure}[htbp]
\centering
\includegraphics[width=\linewidth]{D:/github/Xbotics-Learning-Recording/Xbotics-Learning-Recording/Week3/figures/fig_rl_qlearning_dqn.png}
\caption{Q-learning 与 DQN。(a) 三种配置的训练曲线：两种表格法都卡在 150 分附近，DQN 打满 500；(b) 表格 Q 表学到的策略在 $\theta$–$\dot{\theta}$ 平面上是分块常数；(c) DQN 的 Q 值差（推右减推左）是一张光滑曲面，决策边界连续。}
\end{figure}

\begin{table}[htbp]
\centering
\begin{tabularx}{\textwidth}{>{\hsize=0.9000\hsize\raggedright\arraybackslash}X>{\hsize=0.5000\hsize\raggedright\arraybackslash}X>{\hsize=0.7500\hsize\raggedright\arraybackslash}X>{\hsize=1.8500\hsize\raggedright\arraybackslash}X}
\toprule
\rowcolor{thd}\textcolor{white}{\textbf{配置}} & \textcolor{white}{\textbf{状态数}} & \textcolor{white}{\textbf{贪心回报}} & \textcolor{white}{\textbf{说明}} \\
\midrule
随机策略 & — & 23.1 ± 14.8 & 任务难度下界 \\
\rowcolor{tal}
恒定动作 & — & 9.3 ± 0.7 & 最差参考 \\
表格 Q-learning & 4$^{4}$ = 256 & 139.1 ± 67.3 & 离散化粗，状态覆盖好 \\
\rowcolor{tal}
表格 Q-learning & 8$^{4}$ = 4096 & 152.7 ± 78.9 & 格子更细但更稀疏，提升有限 \\
DQN（最后一个网络） & 连续 & 500.0 ± 0.0 & 本随机种子下也达标 \\
\rowcolor{tal}
DQN（最优快照） & 连续 & 500.0 ± 0.0 & 周期贪心评估选出的最好一版 \\
\bottomrule
\end{tabularx}
\caption{值方法的实测结果（脚本 05）。表格法的天花板是离散化误差，不是调参问题。}
\end{table}

DQN 的训练过程本身也很能说明问题。脚本每 5000 步做一次贪心评估，实测序列为 5k:184 $\rightarrow$ 10k:99 $\rightarrow$ 15k:171 $\rightarrow$ 20k:125 $\rightarrow$ 25k:141 $\rightarrow$ 30k:180 $\rightarrow$ 35k:198 $\rightarrow$ 40k:473 $\rightarrow$ 45k:500 $\rightarrow$ 50k:500 $\rightarrow$ 55k:500（环境步 : 回报）。可以看到 10k 步时从 184 掉到 99 —— \textbf{训练到一半策略崩了}，之后才慢慢爬回来，40k 步后稳定在 500。这解释了为什么工程上一定要做周期评估 + 保留最优快照：只交最后一个网络是不可靠的。\par

\hTwo{4.6 三种策略梯度变体与方差实验}

三种变体在相同的 1000 个 episode 预算下：\par

\begin{table}[htbp]
\centering
\begin{tabularx}{\textwidth}{>{\hsize=1.0462\hsize\raggedright\arraybackslash}X>{\hsize=0.9846\hsize\raggedright\arraybackslash}X>{\hsize=1.0462\hsize\raggedright\arraybackslash}X>{\hsize=0.9231\hsize\raggedright\arraybackslash}X}
\toprule
\rowcolor{thd}\textcolor{white}{\textbf{方法}} & \textcolor{white}{\textbf{贪心回报}} & \textcolor{white}{\textbf{平均 episode 长度}} & \textcolor{white}{\textbf{平均策略熵}} \\
\midrule
REINFORCE & 125.4 ± 7.5 & 371.1 & 0.533 \\
\rowcolor{tal}
REINFORCE + 基线 & 500.0 ± 0.0 & 397.6 & 0.533 \\
A2C & 500.0 ± 0.0 & 322.2 & 0.410 \\
\rowcolor{tal}
随机策略 & 23.1 ± 14.8 & — & — \\
\bottomrule
\end{tabularx}
\caption{三种策略梯度变体的实测结果（脚本 06，1000 个 episode）。REINFORCE 在这个预算下没能解掉任务，加上基线后立刻达标 ——方差就是这个方法能不能工作的分水岭。}
\end{table}

\begin{figure}[htbp]
\centering
\includegraphics[width=\linewidth]{D:/github/Xbotics-Learning-Recording/Xbotics-Learning-Recording/Week3/figures/fig_rl_policy_gradient.png}
\caption{三种策略梯度变体。(a) 锯齿状的回报曲线是策略梯度方法方差大的直观表现 ——冲上去、崩下来、再学回来；(b) 固定策略下，减去基线把信用分配信号的标准差从 8.53 压到 5.74；(c) 策略熵的下降速度。}
\end{figure}

受控方差实验的结果：固定策略、同一批 60 条轨迹，只把优势函数从 $G_{t}$ 换成 $G_{t}$ $-$ V($s_{t}$)，信用分配信号的标准差从 \textbf{8.53 降到 5.74}，其中评论家解释了 54.7\% 的回报方差。这直接验证了周报 4.4 中“减基线不引入偏差、却能压低方差”的结论。\par

\hTwo{4.7 PPO 与裁剪消融}

\begin{figure}[htbp]
\centering
\includegraphics[width=\linewidth]{D:/github/Xbotics-Learning-Recording/Xbotics-Learning-Recording/Week3/figures/fig_rl_ppo.png}
\caption{PPO 与消融。(a) 回报 vs 环境步数：橙线是 PPO 训练时随机策略的回报，绿虚线是周期贪心评估；(b) 每次更新的近似 KL，裁剪后全程贴在 1e-2 附近，不裁剪则在 1e-2 \textasciitilde{} 1e0 之间剧烈震荡；(c) 裁剪真正生效的样本比例，稳定在 10\% 左右。}
\end{figure}

\begin{table}[htbp]
\centering
\begin{tabularx}{\textwidth}{>{\hsize=1.1000\hsize\raggedright\arraybackslash}X>{\hsize=0.9500\hsize\raggedright\arraybackslash}X>{\hsize=0.9500\hsize\raggedright\arraybackslash}X>{\hsize=1.0000\hsize\raggedright\arraybackslash}X}
\toprule
\rowcolor{thd}\textcolor{white}{\textbf{指标}} & \textcolor{white}{\textbf{PPO（$\varepsilon$ = 0.2）}} & \textcolor{white}{\textbf{不裁剪（消融）}} & \textcolor{white}{\textbf{差距}} \\
\midrule
平均每次更新的近似 KL & 0.01449 & 0.26130 & 18.0 倍 \\
\rowcolor{tal}
最大 KL & 0.08523 & 1.03419 & 12.1 倍 \\
裁剪生效的样本比例 & 10.4\% & 0.0\% & — \\
\rowcolor{tal}
贪心回报 & 500.0 & （未单独评估） & — \\
\bottomrule
\end{tabularx}
\caption{裁剪消融实验。两组的学习率、批大小、epoch 数完全相同，唯一差别是裁剪是否开启；不裁剪时每次更新把策略推得远约 18 倍。}
\end{table}

也就是说，“稳定”可以量化成一句：\textbf{裁剪把每次更新的策略变化量（以近似 KL 度量）压小了约 18 倍}，而且这不是靠调小学习率换来的。裁剪生效的样本比例稳定在 10.4\% 左右，说明这项机制确实在持续工作，而不是形同虚设。\par

\hTwo{4.8 A3C 并行采样与横向汇总}

固定策略下，让每个“worker”跑一条轨迹，统计“平均 B 条轨迹的梯度估计”的标准差：\par

\begin{table}[htbp]
\centering
\begin{tabularx}{\textwidth}{>{\hsize=1.2963\hsize\raggedright\arraybackslash}X>{\hsize=0.9506\hsize\raggedright\arraybackslash}X>{\hsize=0.9506\hsize\raggedright\arraybackslash}X>{\hsize=0.9506\hsize\raggedright\arraybackslash}X>{\hsize=0.9506\hsize\raggedright\arraybackslash}X>{\hsize=0.9506\hsize\raggedright\arraybackslash}X>{\hsize=0.9506\hsize\raggedright\arraybackslash}X}
\toprule
\rowcolor{thd}\textcolor{white}{\textbf{并行 worker 数 B}} & \textcolor{white}{\textbf{1}} & \textcolor{white}{\textbf{2}} & \textcolor{white}{\textbf{4}} & \textcolor{white}{\textbf{8}} & \textcolor{white}{\textbf{16}} & \textcolor{white}{\textbf{32}} \\
\midrule
梯度估计标准差 & 0.3411 & 0.2509 & 0.1958 & 0.1464 & 0.0932 & 0.0749 \\
\rowcolor{tal}
相对单 worker & 100.0\% & 73.5\% & 57.4\% & 42.9\% & 27.3\% & 22.0\% \\
\bottomrule
\end{tabularx}
\caption{并行采样对梯度方差的削减。32 个 worker 把标准差压到单 worker 的 22.0\%（理论上限 $1/\sqrt{32}\approx 17.7\%$）。}
\end{table}

对数据做 log-log 幂律拟合，得到的幂指数为 0.447，与理论值 0.5（即标准差按 $1/\sqrt{B}$ 下降、方差按 1/B 下降）吻合得很好，定量验证了周报 4.6 中 A3C 的方差缩减结论。\par

\begin{figure}[htbp]
\centering
\includegraphics[width=\linewidth]{D:/github/Xbotics-Learning-Recording/Xbotics-Learning-Recording/Week3/figures/fig_rl_summary.png}
\caption{横向汇总。(a) 各算法的最终贪心回报；(b) A3C 并行采样的方差缩减，实测幂律指数 0.447（理论 0.5）；(c) 两个“稳定化改造”的量化效果对比。}
\end{figure}

\begin{table}[htbp]
\centering
\begin{tabularx}{\textwidth}{>{\hsize=1.1500\hsize\raggedright\arraybackslash}X>{\hsize=0.6500\hsize\raggedright\arraybackslash}X>{\hsize=0.4750\hsize\raggedright\arraybackslash}X>{\hsize=1.7250\hsize\raggedright\arraybackslash}X}
\toprule
\rowcolor{thd}\textcolor{white}{\textbf{方法}} & \textcolor{white}{\textbf{类别}} & \textcolor{white}{\textbf{贪心回报}} & \textcolor{white}{\textbf{关键机制}} \\
\midrule
随机策略 & 基线 & 23.1 & CartPole 随机动作，任务难度下界 \\
\rowcolor{tal}
恒定动作 & 基线 & 9.3 & 始终推左，最差参考 \\
表格 Q-learning (256 状态) & 值迭代 / 表格 & 139.1 & 离散化粗，状态覆盖好 \\
\rowcolor{tal}
表格 Q-learning (4096 状态) & 值迭代 / 表格 & 152.7 & 格子更细但更稀疏，提升有限 \\
DQN（最优快照） & 值函数逼近 & 500.0 & 连续观测 + 经验回放 + 目标网络 \\
\rowcolor{tal}
DQN（最后一个网络） & 值函数逼近 & 500.0 & 同一配置的末态；训练中会崩，靠快照选择兜底 \\
REINFORCE & 策略梯度 & 125.4 & 无基线，梯度方差最大 \\
\rowcolor{tal}
REINFORCE + 基线 & 策略梯度 & 500.0 & 减去蒙特卡洛评论家 V(s) \\
A2C & 演员-评论家 & 500.0 & 单步自举优势，偏差大方差小 \\
\rowcolor{tal}
PPO (clip=0.2) & 演员-评论家 & 500.0 & 裁剪 + 同批数据复用 10 个 epoch \\
\bottomrule
\end{tabularx}
\caption{十种方法与基线的横向对比（脚本 08 汇总 05/06/07 的结果文件）。}
\end{table}

这张表还带出一个诚实的观察：\textbf{在 CartPole 上，好几种方法最后都能打满 500 分，单看最终分数区分不出算法优劣}。真正的差异在过程里 —— 样本效率、梯度方差、训练曲线的抖动程度、以及对超参和随机种子的敏感度。所以本文才把“信用分配信号的方差”和“每次更新的 KL”这两个过程指标单独量了出来。\par

样本效率方面，PPO 用 10 万环境步达到满分，贪心评估在 81920 步首次摸到 500。作为对照，REINFORCE 在同样的 1000 个 episode 预算下（约 371k 环境步，是 PPO 总预算的 3.7 倍）\textbf{根本没解掉任务}：贪心评估只有 125.4 分，而同一个训练循环一旦加上基线就变成 500 分。这个对照比单看分数更有意义：\textbf{REINFORCE 能不能工作，几乎完全取决于有没有一个好基线}，而 PPO 把这件事做成了算法的一部分。\par

还有一个值得记下的现象：REINFORCE 训练期间采样策略的平均回报能到 303 分，但把策略换成确定性的贪心动作后只剩 125.4 分。说明它的随机策略是“靠探索东一下西一下平均出来的”，而它的贪心路径本身是错的。这类“随机比贪心好”的策略在策略梯度里并不罕见，也是评估时必须固定用贪心、且要单独报告的原因。\par

\clearpage

\hOne{五、问题与解决}

\hTwo{5.1 实现层面踩过的坑}

\begin{table}[htbp]
\centering
\begin{tabularx}{\textwidth}{>{\hsize=0.7875\hsize\raggedright\arraybackslash}X>{\hsize=1.0500\hsize\raggedright\arraybackslash}X>{\hsize=1.1625\hsize\raggedright\arraybackslash}X}
\toprule
\rowcolor{thd}\textcolor{white}{\textbf{问题现象}} & \textcolor{white}{\textbf{根本原因}} & \textcolor{white}{\textbf{解决办法}} \\
\midrule
两连杆正运动学与解析解对不上 & make\_transform(R, p) 构造的是“位姿”[[R, p]]，而连杆变换是 $\mathrm{Rot}_z(\theta)\cdot\mathrm{Trans}(l)$，两者的平移列不同；混用导致整条链错位 & 拆成 rot\_z\_h / trans 两个显式 helper，让复合顺序在代码里一眼可见 \\
\rowcolor{tal}
DH 链检查全部通过，但 FK 结果全错 & 链式检查用的是同一个错误的 dh\_transform，属于“自己验证自己”的假阴性；另外 DH 展开时漏掉 $R_{z}$($\theta$) 对平移列的作用 & 加一条独立检查，把 DH 矩阵与教科书展开式逐项对照 \\
RRT 路径与障碍物重叠 & 线段碰撞检测的采样步长太粗，会从薄障碍中间穿过去 & 把采样步长降到 5 mm（远小于最小障碍尺寸），并留注释说明升级路径 \\
\rowcolor{tal}
表格 Q-learning 稳定卡在 150 分，加密格子反而更差 & 不是 bug，是离散化误差的天花板：格子越密，每个格子见到的样本越少 & 用 DQN（连续观测）做对照确认结论，并把这一现象写进结果说明 \\
DQN 训练到一半策略崩溃（贪心评估 184 $\rightarrow$ 99） & 每步都做梯度更新 + 学习率 1e-3 过于激进，Q 值持续膨胀（|Q|max 从 8 涨到 51） & 学习率降到 2.5e-4，并按固定间隔做贪心评估、保留最优快照 \\
\rowcolor{tal}
REINFORCE 偶尔塌缩到 9 分（比随机还差） & 学习率 1e-2 时策略过早坍缩成确定性动作，探索彻底停止 & 学习率降到 1e-3，同时保留熵奖励项 \\
\bottomrule
\end{tabularx}
\caption{实现层面遇到的主要问题与解决办法。}
\end{table}

最值得记的一条是第二行那个“假阴性”：\textbf{用同一个（错误的）实现去验证它自己，断言再多也没用}。后来我给所有关键数学都配了独立的交叉验证源 ——正运动学对解析闭式解、解析雅可比对有限差分、SLERP 对解析转角、PPO 的 GAE 对 $\lambda$ = 0 与 $\lambda$ = 1 两个退化端点。\par

\hTwo{5.2 实验可复现性}

训练类脚本的可复现性比数值脚本更容易出问题，中途踩了两个坑：\par

\begin{itemize}[leftmargin=1.6em,itemsep=0.15em,topsep=0.25em]

\item \textbf{环境未播种。}只设 torch/numpy 的随机种子是不够的：episode 的初始状态由环境的随机数发生器决定，而 gymnasium 在创建环境时会用系统熵给它播种。结果就是“同样的种子，两次运行的指标不同”。解决办法是在训练开始时显式调用一次 env.reset(seed=seed)，后续的 reset 都从这条已播种的随机流继续推进 ——既保证初始状态多样，又保证整个训练过程可复现。

\item \textbf{随机动作未播种。}更隐蔽的一个：env.reset(seed=...) 只保证初始状态可复现，而 env.\allowbreak action\_\allowbreak space.\allowbreak sample() 用的是另一个随机源，不单独播种就每次都不一样。这会让“随机策略基线”这个参照值在每次运行时漂移，解决办法是显式调用 env.\allowbreak action\_\allowbreak space.\allowbreak seed(...)。

\end{itemize}

两处修完后，重复运行 06 与 08 得到的 JSON 结果文件完全一致（逐字段比对通过），本说明中的所有数字都可复现。\par

另外还有一处与可复现性相关的设计取舍值得说明：DQN 脚本采用了“周期贪心评估 + 保留最优快照”。严格来说，这等于用评估集做了模型选择，所以报告中同时给出“最后一个网络”与“最优快照”两个数，两者之间的差距本身就是 DQN 训练不稳定的证据 ——而不是把它藏起来只报最好看的那个数。\par

\hOne{六、结论}

本套代码对周报中的每一条理论结论都给出了可执行的验证：\par

\begin{itemize}[leftmargin=1.6em,itemsep=0.15em,topsep=0.25em]

\item \textbf{理论结论 $\rightarrow$ 验证方式：}位姿变换的约定 $\rightarrow$ 往返一致性断言；链式法则 $\rightarrow$ 与解析闭式解对照；雅可比 $\rightarrow$ 与有限差分对照；速度/加速度上限 $\rightarrow$ 与闭式峰值公式对照；“两秒可行性判据” $\rightarrow$ 与网格搜索最优解对照；Bellman 目标 $\rightarrow$ 自举掩码断言；裁剪目标的四种情形 $\rightarrow$ 逐项梯度断言；A3C 的 1/B 方差缩减 $\rightarrow$ 固定策略下的受控实验与幂律拟合。

\end{itemize}

三个最值得记住的实测结论：\par

\begin{itemize}[leftmargin=1.6em,itemsep=0.15em,topsep=0.25em]

\item 三次多项式只给端点条件时峰值会超限（$\Delta$q = 1 rad、T = 1 s 时峰值加速度 6 rad/s$^{2}$，是给定的 2 rad/s$^{2}$ 上限的 3 倍）——限制必须显式写进规划器，不能靠“留余量”。

\item 表格 Q-learning 稳定卡在 150 分附近，加密格子反而更差；DQN 直接吃连续观测可以打满 500 ——格子法的瓶颈是离散化误差，不是调参。

\item PPO 的裁剪把每次更新的近似 KL 从 0.261 压到 0.01449（约 18 倍），而学习率、批大小、epoch 数都完全相同 ——“更稳定”是可以量化的，不需要停留在定性描述。

\end{itemize}

\end{document}
